% !TeX program = xelatex
% ============================================================
% research/report/report.tex — Отчёт исследовательского курсового
% проекта ОП «Программная инженерия» ФКН НИУ ВШЭ.
%
% Оформление: ГОСТ 7.32-2017 «Отчёт о НИР». Чистый документ БЕЗ штампа
% (в отличие от ЕСПД-документов прикладного курсового). Структура и язык
% ветвятся по project.lang:
%   RU — форма ОП ПИ (ru pi): Список исполнителей → Реферат → Содержание →
%        Термины → Введение → главы (обзор→метод→результаты, «Выводы по
%        главе») → Заключение → Список источников (ГОСТ Р 7.0.5-2008) →
%        Приложения;
%   EN — англ. академическая форма ФКН: Annotation → Keywords → Contents →
%        Introduction → Literature review → Methodology → Experiments &
%        results → Conclusion → References → Appendices.
%
% Как пользоваться: пишите свой текст вместо жёлтых \hseFill{…}; данные о
% себе, руководителе, группе и годе задайте в настройках проекта.
% ============================================================
\documentclass[12pt,a4paper]{extreport}
\input{preamble}
\begin{document}
\input{title}


% ============================================================
%  ENGLISH RESEARCH REPORT (project.lang == en)
% ============================================================
\chapter*{Annotation}
\addcontentsline{toc}{chapter}{Annotation}
% A concise summary of the study: object, aim, methods, main results and
% their applicability (a short paragraph).
This report presents \hseFill{a concise summary: the problem, the aim, what was done and the main result}. The object of the study is \hseFill{object}; the aim is \hseFill{aim}. \hseFill{Briefly: methods used and the main result obtained, and where it can be applied.}

\clearpage
\chapter*{Keywords}
\addcontentsline{toc}{chapter}{Keywords}
% 5–15 keywords/phrases, comma-separated.
\hseFill{Keyword one, keyword two, keyword three, keyword four, keyword five}

\clearpage
% ── Обязательный русскоязычный реферат ──
% Требование §2.3.4.2 Программы практики ОП ПИ: исследовательская работа,
% выполненная на английском языке, обязана содержать И англо-, И
% русскоязычный реферат (не опционально). Печатается на русском независимо
% от языка проекта.
\begin{otherlanguage}{russian}
\chapter*{Реферат}
\addcontentsline{toc}{chapter}{Реферат}
Отчёт содержит \hseFill{краткое изложение: проблема, цель, что сделано и главный результат}. Объект исследования~--- \hseFill{объект}; цель работы~--- \hseFill{цель}. \hseFill{Кратко: использованные методы, основной полученный результат и область его применения.}

\vspace{1em}
\noindent\textbf{Ключевые слова:} \hseFill{ключевое слово один, ключевое слово два, ключевое слово три, ключевое слово четыре, ключевое слово пять}
\end{otherlanguage}

\clearpage
\tableofcontents
\clearpage

% ============================================================
\chapter{Introduction}
% ============================================================
% Relevance, aim, objectives, object and subject of the study, and the
% structure of the report.
Motivate the topic: what problem is addressed, why it matters, and what is
already known~\cite{example-en}.

\textbf{Aim of the study} — \hseFill{one sentence: the scientific result to be achieved}.

\textbf{Objectives:}
\begin{enumerate}[label=\arabic*), leftmargin=1.25cm]
  \item \hseFill{review existing approaches};
  \item \hseFill{propose / formalise a method (model, algorithm)};
  \item \hseFill{implement it and run a computational experiment};
  \item \hseFill{analyse the results and draw conclusions}.
\end{enumerate}

\textbf{Object of study} — \hseFill{what is studied in general}.
\textbf{Subject of study} — \hseFill{the specific aspect in focus}.

\clearpage
% ============================================================
\chapter{Literature review}
% ============================================================
\section{State of the art}
Give an analytical overview of existing results and approaches, with
citations~\cite{example-en}.

\section{Comparison of existing approaches}
Compare known methods by criteria relevant to the study; the comparison is
given in Table~\ref{tab:approaches-en}.

\begin{table}[H]
\caption{Comparison of existing approaches}
\label{tab:approaches-en}
\centering
\begin{tabular}{|>{\raggedright\arraybackslash}p{5cm}|c|c|c|}
\hline
\textbf{Criterion} & \textbf{Approach~A} & \textbf{Approach~B} & \textbf{Proposed} \\
\hline
Criterion~1 & \cmpplus  & \cmpminus  & \cmpplus \\
\hline
Criterion~2 & \cmpminus & \cmpplanned & \cmpplus \\
\hline
\end{tabular}
\end{table}

\section{Summary}
State the gap in existing approaches and the resulting problem statement.

\clearpage
% ============================================================
\chapter{Methodology}
% ============================================================
\section{Problem statement}
State the studied problem formally: input, expected output, constraints and
the quality criterion (e.g. an objective $f(x)\to\min$ or a hypothesis).

\section{Proposed method}
Describe the proposed method (model, algorithm): its main steps, assumptions
and parameters.

\section{Summary}
State the key methodological results.

\clearpage
% ============================================================
\chapter{Experimental setup and results}
% ============================================================
\section{Experiment design, data and metrics}
Describe the research questions, datasets (size, source, preparation) and the
quality metrics used.

\section{Results and discussion}
Report the results (Table~\ref{tab:results-en}), compare with baselines and
discuss significance and limitations.

\begin{table}[H]
\caption{Experimental results}
\label{tab:results-en}
\centering
\begin{tabular}{|>{\raggedright\arraybackslash}p{5cm}|c|c|}
\hline
\textbf{Method} & \textbf{Metric~1} & \textbf{Metric~2} \\
\hline
Baseline         & \ldots & \ldots \\
\hline
Proposed method  & \ldots & \ldots \\
\hline
\end{tabular}
\end{table}

\section{Summary}
State the outcomes of the experimental study.

\clearpage
% ============================================================
\chapter{Conclusion}
% ============================================================
Summarise the results against each objective, the main outcome, and directions
for further work.

\clearpage
% ── References (GOST R 7.0.5-2008 or an IEEE-style numbered list) ──
\phantomsection
\addcontentsline{toc}{chapter}{References}
\renewcommand{\bibname}{References}
\begin{thebibliography}{99}
  \bibitem{example-en}
  Author A.~B., Coauthor C.~D. Title of the paper // Proceedings of the Conference. --- 2024. --- P.~10--20.
  \bibitem{example-web}
  Resource name [Electronic resource]. --- URL: \url{https://example.org} (accessed: 01.01.2026).
\end{thebibliography}

\clearpage
% ── Appendices (GOST 7.32-2017 §6.14) — optional ──
\appendix
\renewcommand{\thechapter}{\Alph{chapter}}
\titleformat{\chapter}[hang]
  {\normalfont\large\bfseries\centering}{Appendix~\thechapter}{1ex}{}
\chapter{Illustrations}
Place bulky figures, listings and data tables here.



\end{document}
